% ===================================================================
%  કોષ્ટક નં. ૨ — માનવ શરીર. --- વિરામનો સમય. રમતો.
%
%  Traduction, faite en 2026, en gujarati d'aujourd'hui, sur l'ido de
%  texto/io. Regles de la colonne : voir 10-kostak-01.tex. Une seule
%  numerotation. Ecrit parigi.py ouvert, apres les deux fautes du
%  tableau 1.
%
%  LE CORPS A TROIS REGISTRES ICI, ET LA COLONNE DOIT CHOISIR A
%  CHAQUE MOT. Le gujarati garde pour presque chaque partie du corps
%  un mot SANSKRIT PUR (tatsama), un mot USE par les siecles
%  (tadbhava), et souvent un troisieme, PERSAN, venu du Sultanat :
%
%    મસ્તક / માથું          la tete
%    નેત્ર / આંખ            l'oeil
%    ઉદર / પેટ             le ventre
%    હૃદય / દિલ (persan)   le coeur
%    મુખ / મોં             la bouche
%
%  LA REGLE DE 2026 EST NETTE ET LA COLONNE S'Y TIENT : le tadbhava
%  est ce qu'on DIT, le tatsama est ce qu'on ECRIT dans un livre de
%  science. Or l'ido annonce lui-meme cet alinea comme « simpla
%  leciono pri naturcienco » donnee a des enfants. On ecrit donc
%  માથું, આંખ, પેટ, હાથ, પગ — la langue parlee — et l'on garde le
%  sanskrit pour ce que la lecon nomme en savant : les organes et les
%  cinq sens. Le persan દિલ n'est pas ecrit : il vit tres bien en
%  gujarati, mais dans les chansons et les proverbes, pas dans une
%  planche d'anatomie.
%
%  LES CINQ SENS SONT LA QUATRIEME REPONSE DU RELEVE A CET ALINEA, ET
%  ELLE REJOINT LA PERSANE PAR L'AUTRE BOUT. Le persan formait cinq
%  noms sur cinq VERBES PERSANS avec un seul suffixe ; le javanais
%  cinq formes en pa- ; le haoussa n'avait qu'un verbe pour quatre.
%  Le gujarati ecrit દૃષ્ટિ, શ્રવણ, ઘ્રાણ, સ્વાદ, સ્પર્શ : cinq noms
%  d'action sanskrits, batis sur cinq racines verbales par le meme
%  procede — exactement la meme elegance que la persane, mais prise
%  toute faite dans une autre langue au lieu d'etre fabriquee dans la
%  sienne.
%
%  LA SERRE (26) EST « કાચઘર », LA MAISON DE VERRE, ET LES TROIS
%  DERNIERES COLONNES ONT NOMME TROIS CHOSES DIFFERENTES. Le persan
%  disait گلخانه, la maison des FLEURS ; le haoussa « gidan tsirrai »,
%  la maison des PLANTES ; le gujarati la maison du VERRE. Trois
%  langues sans parente devant le meme batiment, et chacune a choisi
%  un element different du meme objet — ce qu'il contient, ce qu'on y
%  fait pousser, ce dont il est fait.
%
%  LES ECHASSES (24) SONT « લાકડાના પગ », LES JAMBES DE BOIS, ET
%  C'EST LA TROISIEME FOIS DE SUITE. Le persan disait چوب‌پا, le
%  bois-pied ; le haoussa « sandunan ƙafa », les batons de jambe.
%  Trois colonnes, trois langues, la meme composition en deux mots.
%
%  PREMIER REFUS DE LA COLONNE : LES QUILLES (16). Le gujarati a
%  « ગિલ્લી-દંડા », le jeu du baton et du fuseau, qui est LA rue de
%  toute enfance gujaratie et qui se joue avec un baton qu'on
%  brandit ; la planche grave neuf quilles alignees qu'on abat d'une
%  boule. On ecrit « લાકડાના દાંડા ».
%
%  MAIS LE CERF-VOLANT (33) EST « પતંગ » SANS UNE SECONDE
%  D'HESITATION, ET C'EST L'OBJET DU LIVRET LE PLUS MAL SERVI PAR SA
%  PLANCHE. Le graveur montre un garcon qui tient une ficelle dans un
%  parc. En Gujarat le patang a un mois, une fete — Uttarayan, le
%  14 janvier —, un ciel entier au-dessus d'Ahmedabad, une bobine qui
%  porte son nom (ફિરકી), un fil enduit de verre pile (માંજા) et un
%  cri quand on coupe celui du voisin (« કાપ્યો છે ! »). Le persan
%  decrivait le sien, بادبادک, le petit vent-vent ; le haoussa
%  l'appelait « tsuntsun iska », l'oiseau de vent. Ici il n'y avait
%  rien a chercher : le mot etait la, et c'est la chose gravee qui
%  est trop petite pour lui.
%
%  ET LA BICYCLETTE (37) EST « સાયકલ », SIXIEME MOT DU RELEVE POUR
%  CETTE MACHINE. pit du neerlandais, sepeda du velocipede,
%  bicyclette, دوچرخه la deux-roues, keke la machine a manivelle — et
%  l'anglais tout cru. Six langues, six solutions, et celle-ci est la
%  seule qui n'ait ni emprunte au voisin, ni decrit, ni reutilise :
%  elle a pris le mot de celui qui vendait l'objet.
% ===================================================================

%%K t02-apar-1 apar
\VUsaut{4.0mm}
\VUtitre{15.0pt}{60}{કોષ્ટક નં. ૨}
\VUsaut{2.0mm}
\VUtitre{11.4pt}{0}{માનવ શરીર. --- વિરામનો સમય.}
\VUtitre{11.4pt}{0}{રમતો.}

%%K t02-tit-0 sub
\VUtitre{13.2pt}{0}{માનવ શરીર.}
\VUsaut{2.4mm}
\VUcentre{10.2pt}{0}{\textit{(પ્રકૃતિવિજ્ઞાનનો સાદો પાઠ.)}}
\VUsaut{3.0mm}

%%K t02-tit-1 sub pos=t02-01-1
\VUpk{10.2pt}{40}{માથું.}
\VUsaut{3.0mm}

%%K t02-01-1 p
1. --- માનવ શરીરના મુખ્ય ભાગ છે : \VUgras{માથું,} \VUgras{ધડ} અને
\VUgras{અંગો.}

%%K t02-02-1 p
2. --- માથામાં બે ભાગ છે : \VUgras{ખોપરી,} જેમાં \VUgras{મગજ} છે અને
જેને \VUgras{વાળ} ઢાંકે છે, અને \VUgras{ચહેરો.}

%%K t02-03-1 p
3. --- અમારા ચિત્રમાં અમે ખોપરી કે મગજ જોતા નથી ; પણ ચહેરાના મહત્ત્વના
ભાગ અમે ઘણી સ્પષ્ટતાથી જોઈએ છીએ.

%%K t02-04-1 p
4. --- પહેલાં આ \VUgras{કપાળ,} જે પ્રબળ બુદ્ધિ અને સતત ચાલતા વિચારને
બતાવે છે. \VUgras{લમણાં} ઘણાં ખુલ્લાં છે. \VUgras{આંખ} જીવંત અને પહોળી
ખૂલેલી છે. ઘાટી \VUgras{ભમર} અને લાંબી \VUgras{પાંપણો} તેને રક્ષે છે.

%%K t02-05-1 p
5. --- વળી આ \VUgras{નાક} અને \VUgras{નસકોરાં.} નાક અમને સૂંઘવા અને
શ્વાસ લેવા કામ આવે છે. નાકની બંને બાજુએ \VUgras{ગાલ} છે, અને વધારે
દૂર \VUgras{કાન.} ચહેરાનો નીચેનો ભાગ \VUgras{મોં} અને \VUgras{દાઢી}નો
બનેલો છે.

%%K t02-06-1 p
6. --- અમે તેમને બરાબર જોતા નથી, કેમ કે \VUgras{મૂછ} અને
\VUgras{ગલમૂછ} તેમને અમારાથી છુપાવે છે. પણ અમે જાણીએ છીએ કે મોંમાં બે
\VUgras{હોઠ,} \VUgras{ઉપલું જડબું} અને \VUgras{નીચલું જડબું} તેમના
\VUgras{દાંત} સાથે, અને \VUgras{જીભ} છે. મોંથી અમે બોલીએ છીએ અને
ખાઈએ છીએ. જીભથી અમે ખોરાકનો સ્વાદ લઈએ છીએ.

%%K t02-07-1 p
7. --- આંખ, કાન, નાક અને મોં, આંગળીઓની જેમ, પાંચ ઇન્દ્રિયોનાં અંગ છે :
\VUgras{દૃષ્ટિ,} \VUgras{શ્રવણ,} \VUgras{ઘ્રાણ,} \VUgras{સ્વાદ,}
\VUgras{સ્પર્શ.}

%%K t02-08-1 p
8. --- માથું આમ માનવ શરીરના સૌથી રસપ્રદ ભાગોમાંનું એક છે.

%%K t02-tit-2 sub
\VUsaut{4.4mm}
\VUpk{10.2pt}{40}{ધડ.}
\VUsaut{3.0mm}

%%K t02-09-1 p
9. --- \VUgras{ગરદન} માથાને ધડ સાથે જોડે છે. તેમાં \VUgras{ગળું} અને
\VUgras{ડોકનો પાછલો ભાગ} છે.

%%K t02-10-1 p
10. --- ધડમાં ઘણાં આવશ્યક અંગ પણ છે. \VUgras{છાતીમાં} \VUgras{ફેફસાં}
છે, જેનાથી અમે શ્વાસ લઈએ છીએ, અને \VUgras{હૃદય,} જે જીવનનું કેન્દ્ર
છે. હૃદય ધબકતું બંધ થાય ત્યારે જીવ મરી જાય છે.

%%K t02-11-1 p
11. --- \VUgras{પાંસળીઓ} છાતીને ટકાવી રાખે છે.

%%K t02-12-1 p
12. --- ધડના બાકીના ભાગ છે \VUgras{કૂખ,} \VUgras{પીઠ,} \VUgras{ખભા}
અને \VUgras{પેટ.} અમે સૂવા માટે કૂખ પર કે પીઠ પર આડા પડીએ છીએ, અમે
ભાર અમારા ખભા પર ઊંચકીએ છીએ.

%%K t02-13-1 p
13. --- પેટમાં \VUgras{હોજરી} અને \VUgras{આંતરડાં} છે. તે અમને ખોરાક
પચાવવામાં કામ આવે છે.

%%K t02-tit-3 sub
\VUsaut{4.4mm}
\VUpk{10.2pt}{40}{અંગો.}
\VUsaut{3.0mm}

%%K t02-14-1 p
14. --- અમારે ચાર \VUgras{અંગ} છે : બે \VUgras{હાથ} અને બે
\VUgras{પગ.} હાથથી અમે લઈએ છીએ, પકડીએ છીએ, ઊંચકીએ છીએ અને વસ્તુઓ ભેગી
કરીએ છીએ ; પગથી અમે ઊભા રહીએ છીએ, ચાલીએ છીએ અને દોડીએ છીએ.

%%K t02-15-1 p
15. --- \VUgras{ખભો} હાથને શરીર સાથે જોડે છે. ઘણો મજબૂત
\VUgras{સ્નાયુ,} એટલે કે બાહુનો સ્નાયુ, હાથને હલાવે છે, જેને
\VUgras{કોણી} \VUgras{આગલા હાથ} સાથે જોડે છે. \VUgras{કાંડું}
\VUgras{પંજાનો} સાંધો બનાવે છે, જે \VUgras{આંગળીઓ} અને \VUgras{નખ}થી
પૂરો થાય છે. હાથ પર અમે \VUgras{નસ} (અને ધમની) ઓળખીએ છીએ, જેમાં
\VUgras{લોહી} વહે છે.

%%K t02-16-1 p
16. --- પગના ભાગ છે : \VUgras{જાંઘ,} \VUgras{ઘૂંટણ} અને \VUgras{ઘૂંટણની
પાછળનો ખાડો,} \VUgras{પિંડી,} જેનું \VUgras{માંસ} \VUgras{ચામડીથી}
ઢંકાયેલું છે, \VUgras{ઘૂંટી} અને \VUgras{પગ.} પગનાં \VUgras{હાડકાં}
ઘણાં જાડાં અને ઘણાં મજબૂત છે. પગના છેડે \VUgras{પગની આંગળીઓ} છે.
બાકીના ભાગ છે \VUgras{પગનું તળિયું} અને \VUgras{એડી.}

%%K t02-17-1 p
17. --- માનવ શરીરનું લગભગ સંપૂર્ણ વર્ણન આવું છે.

%%K t02-17-2 p
\textit{નોંધ.} --- \textit{« મને... (માથામાં વગેરે) દુખે છે » એ
પ્રયોગની મદદથી શિક્ષક આ આખો શબ્દભંડોળ સારી કે ખરાબ}
\VUgras{શારીરિક સ્થિતિ}\textit{ની દૃષ્ટિએ ફરી વાપરી શકશે.}

%%K t02-tit-4 sub
\VUsaut{5.0mm}
\VUpk{10.2pt}{40}{કસરત.}
\VUsaut{2.4mm}
\VUcentre{10.2pt}{0}{\textit{(એક વિદ્યાર્થીએ કહેલી વાત.)}}
\VUsaut{3.0mm}

%%K t02-18-1 p
18. --- લવચીક, ચપળ અને તંદુરસ્ત રહેવા માટે અમારે અમારા પગ અને અમારા
હાથની કસરત કરવી જોઈએ. તેથી અમે રમીએ છીએ અને \VUgras{કસરત} કરીએ છીએ.

%%K t02-19-1 p
19. --- અઠવાડિયા દરમિયાન આ બે કસરતો માધ્યમિક શાળાના મેદાનમાં થાય છે
(કોષ્ટક નં. ૧ જુઓ). પણ ગુરુવારે અને રવિવારે અમે મોટા ઘાસના મેદાનમાં
જઈએ છીએ, જ્યાં અમને દોડવાની અને આનંદ કરવાની જગ્યા મળે છે.

%%K t02-20-1 p
20. --- અમારા શિક્ષકો અને અમારાં માતાપિતા ક્યારેક અમારી સાથે ત્યાં
આવે છે ; પણ કસરતો તો \VUgras{કસરતના શિક્ષક}\textsuperscript{(12)} જ
ચલાવે છે.

%%K t02-21-1 p
21. --- ઘાસના મેદાન પર, પ્રવેશદ્વારની ડાબી બાજુએ,
\VUgras{કસરતનાં સાધનો}\textsuperscript{(1)} છે ; અમે
\VUgras{નિસરણી}\textsuperscript{(2)} પર ચઢીએ છીએ,
\VUgras{કડાં}\textsuperscript{(3)} અને \VUgras{ઝૂલતા
દાંડા}\textsuperscript{(4)} પર ઊંધા લટકીએ છીએ, અમે હાથે હાથે
\VUgras{દોરીની નિસરણી}\textsuperscript{(5)} કે \VUgras{ગાંઠવાળી
દોરી}\textsuperscript{(6)} ની ટોચ સુધી ચઢીએ છીએ.

%%K t02-22-1 p
22. --- એ દરમિયાન અમારા સાથીઓ \VUgras{લાકડાના
ઘોડા}\textsuperscript{(7)} કે \VUgras{સમાંતર દાંડા}\textsuperscript{(11)}
પર કૂદે છે. બીજા \VUgras{મુક્કાબાજી}\textsuperscript{(8)} કરે છે કે
મહોરું અને \VUgras{પાતળી તલવાર} લઈને
\VUgras{તલવારબાજી}\textsuperscript{(9)} કરે છે. છેલ્લે, બીજા
\VUgras{વજન ઉપાડે છે}\textsuperscript{(10)}.

%%K t02-tit-5 sub
\VUsaut{4.6mm}
\VUpk{10.2pt}{40}{રમતો.}
\VUsaut{3.0mm}

%%K t02-23-1 p
23. --- કસરત કરનારાઓની આસપાસ, છોકરા અને છોકરીઓ જુદી જુદી
\VUgras{રમતો} રમે છે. છોકરીઓ પોતાના \VUgras{ગોળ
ચક્ર}\textsuperscript{(14)} થી આનંદ કરે છે ; તેઓ
\VUgras{દોરી}\textsuperscript{(13)} કૂદે છે અથવા પોતાના નાના ભાઈઓ અને
પિતરાઈઓને \VUgras{ક્રોકે}\textsuperscript{(15)} ની રમતમાં બોલાવે છે.
પ્રતિસ્પર્ધીનો \VUgras{દડો} પોતાની \VUgras{લાકડાની નાની હથોડીથી}
ધકેલી દેવામાં તેમને મજા પડે છે, બરાબર તે ઘડીએ જ્યારે તે સહેલાઈથી
કમાનમાંથી પસાર થવાનું વિચારતો હોય !

%%K t02-24-1 p
24. --- છોકરાઓને વધારે તાકાતવાળી રમતો ગમે છે. તેઓ ખુશીથી
\VUgras{લાકડાના દાંડા}\textsuperscript{(16)} ની રમત રમે છે, જેને તેઓ
\VUgras{દડા}\textsuperscript{(17)} વડે કુશળતાથી પાડે છે. તેઓ
\VUgras{ભમરડો}\textsuperscript{(18)} અને \VUgras{દોરીવાળો દડો ને
દાંડી}\textsuperscript{(19)} કે \VUgras{દેડકાની
રમત}\textsuperscript{(20)} પણ ધિક્કારતા નથી. પણ તેમની સૌથી પ્રિય રમતો
તો \VUgras{લખોટી}\textsuperscript{(21)} અને \VUgras{ભીંતદડો}\textsuperscript{(22)}
છે, કેમ કે \VUgras{કાચઘર}\textsuperscript{(26)} ની ભીંત હલકા દડાને
પાછો ઉછાળવા માટે ઘણી અનુકૂળ છે.

%%K t02-25-1 p
25. --- નાનકડો યોહાનેસ, પોતાના \VUgras{લાકડાના
પગ}\textsuperscript{(24)} પર ચાલતો, ઘણો કુશળ છે અને સમતોલન ઘણી સારી
રીતે જાળવી શકે છે. તેની પાસે કેટલાક સાથીઓ તે કાચઘરમાં અને
\VUgras{વાડ}ની પાછળ \VUgras{સંતાકૂકડી}\textsuperscript{(25)} રમે છે.
બીજા \VUgras{થપ્પો}\textsuperscript{(28)} કે
\VUgras{બેડમિન્ટન}\textsuperscript{(29)} વધારે પસંદ કરે છે, જેમાં ફૂલ
એક \VUgras{રેકેટ}થી બીજા રેકેટ સુધી ઊડે છે.

%%K t02-noto-1 noto
\VUnotes{143.2mm}{%
(*) બાળકોની રમતોનાં નામ ઘણી વાર તદ્દન રાષ્ટ્રીય હોય છે. અમે તેમને
ઇડોમાં બને તેટલી સારી રીતે નામ આપવાનો પ્રયત્ન કર્યો છે — ક્યારેક તેમને
ઓળખાવતી ક્રિયામાંથી પ્રેરણા લઈને, ક્યારેક આ કે તે દેશનું રાષ્ટ્રીય
નામ પસંદ કરીને, જ્યારે તે બહુ અયોગ્ય ન હોય. દા. ત. :
\VUgras{પીપડાની રમત} (જર્મન અને ફ્રેન્ચ) એ જ \VUgras{દેડકાની રમત}
\textsuperscript{(20)} (અંગ્રેજી) છે, જેમાં રમનારા પોતાની ગોળ
તકતીઓ પેટીના (પીપડાના) ઉપર બેઠેલા ધાતુના દેડકાના ખુલ્લા મોંમાં
નાખવાનો પ્રયત્ન કરે છે.
\VUgras{ખભાકૂદ}\textsuperscript{(30)} \VUgras{પીઠકૂદ}થી માત્ર આટલામાં
જ જુદો પડે છે : માથા ઉપરથી કૂદવા માટે રમનારા પોતાના હાથ સાથીના ખભા પર
ટેકવે છે, જ્યારે « \VUgras{પીઠકૂદ} » માં તેઓ હાથ માત્ર તેની પીઠ પર જ
ટેકવે છે. --- અથવા તો કેટલાક ભાગ લેનારા વાંકા વળે છે જ્યારે બીજા તેમની
પીઠ ઉપરથી ખભા પર હાથ ટેકવીને કૂદશે ; પહેલાંવાળાએ વળ્યા વિના આઘાત સહન
કરવો પડે, બીજાઓએ સમતોલન ગુમાવ્યા વિના કૂદવું પડે (કોષ્ટક પોતે જુઓ).%
}

%%K t02-26-1 p
26. --- ઘાસના મેદાનની વચ્ચે બે ઝાડ છે. તેમાંના એકના થડ પાસે, સાત કે
આઠ છોકરાઓએ \VUgras{ખભાકૂદ}\textsuperscript{(30)} ની રમત ગોઠવી છે (*).
આ રમત બધા દેશોમાં જાણીતી નથી ; પણ \VUgras{પીઠકૂદ} શું છે તે બધા જાણે
છે, જે આ વખતે છોકરાઓ રમતા નથી.

%%K t02-tit-6 sub
\VUsaut{5.0mm}
\VUpk{10.2pt}{40}{ખુલ્લી હવામાં રમતો.}
\VUsaut{3.4mm}

%%K t02-27-1 p
27. \VUgras{આંધળો પાટો}\textsuperscript{(31)} પણ ઘણી મજાની રમત છે ;
છોકરીઓ અને છોકરાઓ વારાફરતી આંખે \VUgras{પાટો} બાંધે છે અને જે અણઘડ
પકડાઈ જાય તેમને પકડે છે.

%%K t02-28-1 p
28. --- ઘાસના મેદાનની વચ્ચે, આ રહી એક તદ્દન ફ્રેન્ચ પણ થોડી આકરી રમત :
તે છે \VUgras{રીંછની રમત}\textsuperscript{(32)}, જે
\VUgras{પતંગ}\textsuperscript{(33)} ની એટલી શાંત રમત કરતાં કે અંગ્રેજી
રમત \VUgras{ટેનિસ}\textsuperscript{(34)} કરતાં પણ ઘણી વધારે ઘોંઘાટવાળી
છે.

%%K t02-29-1 p
29. --- આ છેલ્લી રમત મોટા વિદ્યાર્થીઓનો આનંદ છે. અમારા શિક્ષકો પોતે
પણ ખુશીથી તે રમે છે. તેમાં ચપળતા સાથે ઘણી કુશળતા જોઈએ છે અને મને તે
ઘણી ગમે છે. \VUgras{હીંચકા}\textsuperscript{(35)} ની રમતો હું છોકરીઓ
માટે છોડી દઉં છું.

%%K t02-30-1 p
30. --- મને બીજી એક અંગ્રેજી રમત ગમતી નથી, જે કદાચ જોખમી બની જાય.

%%K t02-30-1b p
મારો કહેવાનો અર્થ છે \VUgras{ફૂટબોલ}\textsuperscript{(36)}, જે મારા
મોટા ભાઈને આનંદ આપે છે. હું અમારી જૂની \VUgras{દોડપકડ}\textsuperscript{(38)}
ને વધારે પસંદ કરું છું, જે એટલી જ આરોગ્યપ્રદ છે અને ખરેખર ઓછી ક્રૂર
છે.

%%K t02-tit-7 sub
\VUsaut{4.6mm}
\VUpk{10.2pt}{40}{સાયકલ અને દડાની રમત.}
\VUsaut{3.0mm}

%%K t02-31-1 p
31. --- હું ખુશીથી \VUgras{સાયકલ}\textsuperscript{(37)} પર ઘાસના
મેદાનની આસપાસ ફરું છું પણ.

%%K t02-32-1 p
32. --- આવતા વર્ષે હું \VUgras{ઘોડેસવારી}\textsuperscript{(39)} શીખીશ.
જે ઘોડો સુંદર રીતે ડગ ભરે કે દોડે, અને જે કૂદકે અવરોધો ઓળંગે, તે કેટલો
સુંદર લાગે છે !

%%K t02-33-1 p
33. --- ઉનાળામાં અમે \VUgras{તરવાનું}\textsuperscript{(40)} શીખીએ છીએ.
અમે નદીના સ્વચ્છ અને ઠંડા પાણીમાં આનંદથી ડૂબકી મારીએ છીએ અને નહાઈએ
છીએ. ક્યારેક છેલ્લે અમે કાગળનો \VUgras{ફુગ્ગો}\textsuperscript{(41)}
છોડીએ છીએ, જે ગરમ હવાથી ભરાતાં જ ભવ્ય રીતે આકાશમાં ચઢે છે.

%%K t02-34-1 p
34. --- બીજી પણ ઘણી રમતો છે, પણ હું તેમની ગણતરી પૂરી કરી શકું નહીં,
અને આ ટૂંક સારથી સમજાય છે કે ગુરુવારે અમારા સુંદર ઘાસના મેદાન પર બહુ
કંટાળો આવતો નથી.

%%K t02-34-2 p
નોંધ. --- \textit{શિક્ષક આ પ્રકરણ સહેલાઈથી પૂરું કરી શકશે ; સૌથી
મહત્ત્વની રમતોમાંની દરેકનું વધારે વિગતે વર્ણન કરીને.}

\VUsaut{6.0mm}
\VUfilet{22mm}
